\section{Dependency topology and separators}
\label{sec:topology}

Construct a primal dependency graph whose vertices are scalar variables and in
which two variables are adjacent when they occur together in an input
constraint.  A coarse tree decomposition is obtained from the bags
\begin{equation}
  B_i=\{\theta_1,\ldots,\theta_P\}\cup S^{(i)}
  \qquad(i\in[N]),
  \label{eq:sample-bags}
\end{equation}
connected through a central parameter bag
\begin{equation}
  B_0=\{\theta_1,\ldots,\theta_P\}.
\end{equation}

\begin{proposition}[Star decomposition]
\label{prop:star-decomposition}
The bags \(B_0,B_1,\ldots,B_N\), with every \(B_i\) adjacent to \(B_0\), form
a tree decomposition of the coarse sample-local dependency graph.
\end{proposition}

\begin{proof}
Every constraint belongs to one sample and therefore all of its variables lie
in the corresponding bag \(B_i\).  Each local variable occurs in only that bag.
Each parameter occurs in the center and in every sample bag, which is a
connected subtree.  Hence the edge-cover and running-intersection conditions
hold.
\end{proof}

\begin{corollary}[Width bound]
\label{cor:treewidth}
If \(K=\max_iK_i\), then
\begin{equation}
  \boxed{\tw(G)\le P+K-1.}
\end{equation}
\end{corollary}

\begin{proof}
The largest bag in \cref{eq:sample-bags} has at most \(P+K\) vertices.
\end{proof}

Increasing \(N\) adds bags but does not enlarge their shared intersection.
This fact explains why sample-wise linear algebra can scale linearly with
\(N\) once the parameter and local dimensions are fixed.

\begin{remark}[What treewidth does not prove]
A tree decomposition specifies where elimination may occur.  An exact dynamic
program also needs a closed message language whose size and update time are
controlled.  For discrete variables of bounded domain, a table can have size
exponential in the bag width and independent of the total number of bags.  For
real polynomial constraints, a projected semialgebraic set can require new
polynomials, sign conditions, and components.  The numerical dimension of a
separator does not by itself bound this descriptive growth.
\end{remark}

\begin{remark}[Uniform and parameterized readings]
The width bound supports a parameterized analysis in \(P+K\).  A claim of
uniform polynomial time additionally requires polynomial dependence on
\(P+K\), not merely tractability for each fixed width.  The cubic linear-algebra
bound later in the paper has that dependence for evaluated matrices; the exact
semialgebraic message problem does not yet have such a bound.
\end{remark}
